\section{强归纳与良序原理}
\label{sec:01-Mathematical-Language/04-Induction-and-Recursion/02}

\subsection{普通归纳法在这一步哑火了}

回忆一下刚学会的普通归纳法。要证 \(P(n)\) 对所有自然数成立，分两步：基例 \(P(0)\)，然后从 \(P(k)\) 推 \(P(k+1)\)。链条一节一节往前传。

现在来试试证明这样一个命题：Fibonacci 数列的每一项 \(F_n\) 都满足 \(F_n<2^n\)。

Fibonacci 数列的定义是 \(F_0=0,F_1=1\)，而 \(F_n=F_{n-1}+F_{n-2}\) 当 \(n\ge 2\)。基例好办：\(F_0=0<2^0=1\)，\(F_1=1<2^1=2\)。都成立。

现在到了归纳步。假设 \(P(k)\) 成立，即 \(F_k<2^k\)。要证 \(F_{k+1}<2^{k+1}\)。

按定义：\(F_{k+1}=F_k+F_{k-1}\)。我们只知道 \(F_k<2^k\)，但 \(F_{k-1}\) 呢？归纳假设里没有它——我们只假设了 \(P(k)\)，而 \(F_{k-1}\) 对应的是 \(P(k-1)\)。

换句话说，\(F_{k+1}\) 的成立依赖的是前两项，不是前一项。普通归纳法的手，伸不到 \(k-1\) 那个位置去。

这件事不只是在 Fibonacci 上出现。素数分解也类似——把一个合数 \(n\) 拆成 \(a\times b\) 的时候，\(a\) 和 \(b\) 都严格小于 \(n\)，但不一定恰好是 \(n-1\)。普通归纳的``只假设前一个''在这种依赖更远历史的结构面前，有力使不出。

\subsection{既然需要更多假设，那就全拿过来}

强归纳法解决的就是这个问题。它的归纳假设从``只假设 \(P(k)\)''升级为：

\[
\text{假设 }P(n_0),P(n_0+1),\dots,P(k)\text{ 全部成立。}
\]

也就是说，轮到证 \(P(k+1)\) 的时候，你手里可以同时用前面所有已经证明过的命题，想用哪个用哪个。

这听起来比普通归纳法``更强''——假设更多了嘛。但从逻辑的角度看，强归纳和普通归纳能证明的自然数命题集合是完全一样的。为什么？把``从 \(P(n_0)\) 到 \(P(k)\) 都成立''打包成一个新命题 \(Q(k)\)，对 \(Q\) 做一次普通归纳，就能逐字复现强归纳的整个推导过程。强归纳不是拓宽了可证命题的范围，而是让证明的写法更自然地贴合那些``下一步依赖不止一个历史状态''的对象。

\subsection{素数分解：一个非得靠强归纳才能写顺的证明}

\begin{center}
\textbf{命题}：每个大于 1 的整数都可以写成素数的乘积。
\end{center}

拿几个具体数字试试看。2 是素数，行。3 是素数，行。4 不是素数，但 \(4=2\times 2\)，每个因子都是素数。5 是素数。6 不是素数，\(6=2\times 3\)，因子是素数。继续往下试，每遇到一个合数就往更小的因子身上推——但你得先确定那些更小的因子已经能分解。

这就是强归纳法的用武之地。证明对 \(n\ge 2\) 用强归纳。

\textbf{基例} \(n=2\)：2 自身就是素数，当然是素数乘积（一个因子的乘积）。

\textbf{归纳步}：假设从 2 到 \(k\) 的每一个整数都能写成素数乘积。现在看 \(k+1\)：

\begin{itemize}
\tightlist
\item
  如果 \(k+1\) 自己是素数，那它就是素数乘积（一项）。
\item
  如果 \(k+1\) 是合数，那么存在整数 \(a,b\) 满足 \(2\le a,b\le k\) 且 \(k+1=ab\)。由于 \(a\) 和 \(b\) 都在归纳假设的覆盖范围内，它们各自都能写成素数乘积。把这两个素数乘积相乘，\(k+1\) 也就写成了素数乘积。
\end{itemize}

注意归纳假设的用法：我们需要的 \(a\) 和 \(b\) 只是``比 \(k+1\) 小''，并不知道它们具体是多少，也不保证它们刚好等于 \(k\)。强归纳的``全部更小的数都成立''正好兜住了这种不确定性——不管 \(a,b\) 落在 2 到 \(k\) 之间的哪个位置，假设都已经覆盖了。

如果只用普通归纳，就会卡在``\(k+1\) 的两个因子不一定恰好是 \(k\)''这一步上。这就不是证不出来，而是工具不对口。

同时要说清楚：这段证明只完成了一半——证明了分解的存在性，没有碰唯一性。每个整数是不是只有一种方式写成素数乘积（不计因子顺序），还需要算术基本定理的另外一半论证，不能从存在性直接升格。

\subsection{最小反例：从反面攻过来}

换一种思维。不要沿着从小到大的方向推，反过来想：如果命题是假的，那就一定有一个最小的反例——第一个不满足条件的整数。然后证明这个``最早坏掉的''其实不可能坏掉。

这就是\textbf{良序原理}提供的证明策略。良序原理说的是：自然数的任何一个非空子集，都有最小元素。

再用素数分解演示一遍。假设命题不成立，即存在大于 1 的整数不能分解成素数乘积。令 \(n\) 是所有这类数中最小的那个。

\(n\) 不能是素数，因为素数本身已经是素数乘积。所以 \(n\) 是合数：\(n=ab\)，其中 \(2\le a,b<n\)。

关键来了——\(a\) 和 \(b\) 比 \(n\) 小，而 \(n\) 是\emph{最小的}反例。这意味着 \(a\) 和 \(b\) 都不是反例，它们一定可以分解成素数乘积。把 \(a\) 和 \(b\) 的分解乘起来，\(n\) 也就分解了——但这跟 \(n\) 是反例矛盾。

所以反例集合只能是空的。命题成立。

最小反例和强归纳其实是同一枚硬币的两面。强归纳的视角是``如果所有更小的都对，当前也对''；最小反例的视角是``如果当前是第一个出错的，那它更小的部分都没出错，但这就意味着当前也不该出错''。两条路谁先到终点，取决于你的思维习惯——有人觉得正向推导更顺，有人觉得设一个假想敌然后击毙它更爽快。

\subsection{这套工具不能乱搬}

正实数集合就没有最小元素——随便拿一个 \(x>0\)，\(x/2\) 比它更小但还是正数。良序原理不是对任何数集都有效的，它专属于自然数这种离散的、每一步之间没有无限致密空隙的结构。不要看见``取最小值''三个字就条件反射地搬到实数上。

即使用在整数上，``取最小反例''也有一个隐含前提——你得先确认反例集合确实非空。如果命题本来就在所有整数上成立，那反例集合为空，良序原理用不了，论证还没开始就该停了。

更一般地说，要使用最小反例法，需要所讨论的量能嵌入一个良基顺序——不存在无限严格下降的链。对整数的绝对值、有限树的高度、表达式的长度这类度量，通常可以先把问题翻译成自然数上的问题，再调用良序原理。翻译这一步本身也是需要验证的，不是写一句``不妨设''就能带过去。

\subsection{递推结构自然适配强归纳}

回到本节开头那个 Fibonacci 例子。要证 \(F_n<2^n\)，基例 \(n=0,1\) 已检查。对 \(n\ge 2\)，强归纳假设对所有 \(i<n\) 都成立：

\[
F_n = F_{n-1}+F_{n-2} < 2^{n-1}+2^{n-2} = 2^{n-2}(2+1) = 3\cdot 2^{n-2} < 4\cdot 2^{n-2} = 2^n.
\]

前两项都在归纳假设的射程内，一步到位。

凡是递推式回望超过一个历史状态的——二阶梯推、分治算法递归调用、动态规划的状态转移——心里大概就该亮起强归纳的灯了。算法里``所有更小的子问题都已经返回了正确结果''这个前提，正是用强归纳包装的归纳假设。下一节还会把这个模式从整数搬到树和表达式这些非整数结构上。
